Data scarcity remains the primary bottleneck for applying deep learning to experimental photophysics. In this work, we showed that hybrid architectures can improve performance in this setting. By replacing the standard neural head of a trained D-MPNN with a post-hoc XGBoost model, we achieved relatively consistent improvements in absorption wavelength prediction across low-to-mid data regimes ($N \lesssim 1000$), particularly under rigorous scaffold splits that simulate realistic discovery scenarios.

\paragraph{Mechanisms of Stability and Interpretability.}
Our analysis identifies two distinct mechanisms driving this advantage. First, variance decomposition revealed that in small-data regimes, the mapping from learned representations to molecular properties is underdetermined; the hybrid model appears to stabilize this mapping by imposing the structural regularization of decision trees, effectively pruning the high-variance error component associated with neural heads. Second, we showed that the D-MPNN encoder, even when trained on limited data, learns a latent space that aligns with physically meaningful structure. The dominant principal component ($PC_0$) spontaneously aligns with the concept of ``effective conjugation length,'' separating extended $\pi$-conjugation from substituent-driven electronic effects. This suggests that the model’s predictive performance is driven in part by representations aligned with known physical factors, rather than relying solely on dataset-specific correlations.

\paragraph{Practical Implications.}
These findings suggest a practical workflow for molecular screening in low-data settings. Since fitting an XGBoost head on fixed embeddings incurs virtually negligible computational cost, practitioners can leverage the hybrid approach to maximize predictive stability in data-scarce regimes without sacrificing the scalability of end-to-end learning as datasets grow. Furthermore, our transfer learning results indicate that this strategy is particularly effective when adapting pretrained models to chemically similar downstream tasks (e.g., Deep4Chem $\to$ ChemFluor), though its utility diminishes under severe domain shift (e.g., Deep4Chem $\to$ DSSCDB).

\paragraph{Limitations and Future Directions.}
Several open challenges point toward future research.
\begin{itemize}
    \item \textbf{Excited-State Physics:} The hybrid advantage was less systematic for emission prediction. This likely reflects the fact that ground-state graphs do not explicitly encode excited-state relaxation or solvent reorganization energies. Future work should explore integrating low-cost conformational sampling or solvent-aware message passing to better constrain these dynamic effects.
    \item \textbf{Domain Adaptation:} While transfer learning succeeded for structurally overlapping datasets, it failed under significant distribution shift. Addressing this will likely require more principled domain adaptation techniques under data scarcity, such as adversarial pretraining or uncertainty-weighted fine-tuning, rather than simple parameter transfer.
    \item \textbf{Feature Integration:} An additional direction is to evaluate the integration of learned embeddings with physicochemical descriptors within the hybrid framework. Tree-based models naturally accommodate heterogeneous feature sets, suggesting that hybrid approaches may provide a practical mechanism for incorporating complementary domain knowledge in low-data regimes. A controlled comparison with end-to-end models trained on the same augmented inputs is required to determine whether this flexibility translates into improved predictive performance.
    \item \textbf{Scope of Properties:} Extending this analysis to intensive properties like quantum yield or fluorescence lifetime, which depend heavily on non-radiative decay pathways, will further clarify the boundaries of graph-based representation learning.
\end{itemize}

Ultimately, this work supports the use of the hybrid D-MPNN$\rightarrow$XGBoost framework as a strong and practical baseline for optical property prediction in low-data regimes. By combining the representational power of graph neural networks with the robust inference of classical machine learning, it provides a scalable and computationally efficient approach that remains grounded in chemically meaningful structure.